\label{sec:california}

\subsection{Background: The California CRC and Its Criteria}

California's Citizens Redistricting Commission was established by Proposition
11 (2008), extended to congressional redistricting by Proposition 20 (2010).
Unlike the AIRC's five-member structure, the CRC has 14 members: five
registered Democrats, five registered Republicans, and four members registered
with neither major party.

California's redistricting criteria, in statutory priority order, are:
(1) population equality (one person, one vote); (2) VRA compliance; (3) geographic
contiguity; (4) minimization of city and county splits; (5) geographic
compactness; and (6) preservation of communities of interest.  California
does not include partisan fairness as a statutory criterion, distinguishing it
from Arizona.

\subsection{Algorithmic Satisfaction of California's Criteria}

The \texttt{redist} system's weight parameters map onto California's criteria
in a natural way.

\textbf{Criterion 1 (population equality)} is enforced directly: the
algorithm requires population balance within $\pm 0.5\%$ of the
equal-population target, stricter than the legal requirement for congressional
districts.

\textbf{Criterion 2 (VRA compliance)} is addressed through VRASection, which
increases the edge weight on adjacency edges connecting tracts with high
minority-population fractions, encouraging the algorithm to keep minority
communities together.  \citet{dellaLibera2026vra} demonstrates that VRASection
produces more majority-minority districts than the compactness-only baseline.

\textbf{Criterion 3 (geographic contiguity)} is an invariant enforced by the
algorithm's adjacency-graph construction: only geographically contiguous
tracts are considered adjacent, so every output district is contiguous by
construction.

\textbf{Criterion 4 (minimization of city and county splits)} corresponds to
the county-sticky edge weight in the \texttt{redist} system.  Increasing the
weight on county-boundary adjacency edges penalizes splits.  The county-sticky
weight can be calibrated to California's specific county structure, with
larger counties (Los Angeles, San Diego) receiving higher weights to reflect
their larger legislative significance.

\textbf{Criterion 5 (geographic compactness)} is the algorithm's primary
objective through METIS edge-cut minimization.  Polsby-Popper and Reock
compactness scores for California algorithmic districts exceed the CRC's
enacted map on average.

\textbf{Criterion 6 (preservation of communities of interest)} is addressed
through the Commission's adjustment authority, as described for the AIRC in
Section~2.

\subsection{The Proportionality Consideration}

California's criteria do not include partisan proportionality as a statutory
criterion, but CRC commissioners have historically considered proportionality
as a non-binding guideline when evaluating competing plans.  Algorithmic maps
for California tend to produce partisan outcomes close to the statewide
Democratic vote share (approximately 63\% in 2020 presidential results)
because they optimize geometry rather than partisan outcomes.

The relevant question is whether the CRC would find an algorithmic starting
plan adequate.  Given that the algorithm satisfies criteria 1--5 quantifiably
and reproducibly, and given that criterion 6 is reserved for Commission
judgment, the CRC would be in the same position as the AIRC: using the
algorithm output as a constrained starting plan and applying community-of-interest
adjustments where evidence supports them.

\subsection{VRASection and California's Minority Communities}

California has substantial Latino, Asian American, and Black populations that
generate VRA obligations.  The Los Angeles Basin contains multiple majority-Latino
census tracts; the San Francisco Bay Area contains significant Asian American
communities.

VRASection addresses these concentrations by applying enhanced edge weights
that reflect minority population fractions.  The specific threshold for VRA
protection can be set to reflect the CVAP (Citizen Voting Age Population)
percentages required by VRA analysis rather than the total population.
\citet{dellaLibera2026vra} and \citet{dellaLibera2026threshold} develop the
calibration methodology.

Under California's criteria ordering, VRA compliance ranks second (after
population equality), so the VRASection weight is applied before the
county-sticky weight in the algorithm's optimization.  This ensures that
minority-opportunity districts are preserved even where county preservation
would require splitting them.
